This thesis presents TF-Rehancer, a time-frequency speech enhancement architecture for direct full-band processing. The model treats high-frequency components in 48 kHz speech as observed but noisy spectral evidence rather than as missing information to be generated from a band-limited input. TF-Rehancer performs low-band-guided full-band refinement: it encodes low-band speech structure, preserves observed upper-band evidence in the decoder path, and estimates the enhanced spectrum through local complex time-frequency filtering over the original noisy STFT.

We evaluate TF-Rehancer from three perspectives. First, on the 16 kHz VoiceBank+\allowbreak DEMAND benchmark, the compact FastGRU-based TF-Rehancer configuration achieves PESQ 3.33, DNSMOS SIG 3.41, DNSMOS BAK 4.04, and DNSMOS OVL 3.13 with 0.47M parameters and 1.89G model-side MACs. Against the paper-reported FastEnhancer-M reference, it reports higher PESQ, SI-SDR, and DNSMOS P.835 values with fewer model-side MACs, while FastEnhancer-M remains slightly stronger in SCOREQ and ESTOI.

Second, on the VoiceBank+\allowbreak DEMAND 48 kHz test set, the DNS4-trained online TF-Rehancer configuration obtains competitive full-reference scores relative to published PercepNet and DeepFilterNet2 references. Third, controlled 48 kHz ablations show that replacing complex TF filtering with direct spectrum mapping degrades full-reference, spectral, and non-intrusive metrics. They also show that using decoder cross-attention is more effective than removing it on most reported full-reference and spectral metrics, and that a 5 kHz low-band cutoff provides a balanced operating point under the controlled ablation setting.
